\section{Conclusion}
\label{sec:conclusion}

We have given a horizon-wide safety guarantee for closed-loop process operation, obtained by
lifting the per-cycle composed coverage certificate of \citet{busico_m4} with adaptive conformal
inference \citep{gibbs2021} and testing it on a dynamically realised debutanizer loop with
delayed measurements. The guarantee is distribution-free and holds over an operating campaign
rather than at a single cycle.

Four results support it. The conformal health-constrained loop holds the joint safety event at
its certified floor over the horizon, with a $95\%$ Wilson lower bound of $0.992$ against a floor
of $0.75$. A health-blind optimiser meets the quality specification on essentially every cycle
yet over-refluxes until projected life collapses, dropping the joint event to $0.008$, which
shows that meeting a specification and operating safely are distinct and that the binding health
budget is what closes the gap. Validity is invariant to measurement dead time, the selection
penalty staying zero under any label delay, so dead time enters the coverage rate rather than
breaking the guarantee. And adaptive calibration retains a finite interval at target coverage
across horizons at which the naive union bound is vacuous, so adaptivity is necessary rather than
optional.

Taken together these carry the IPIS stack from a certified collection of components to a
certified operating campaign: the per-cycle certificate becomes a runtime, horizon-wide monitor.
The natural next steps are to allocate the similarity budget across several units, assembling a
plantwide certificate from per-unit ones under a shared risk allocation, and to confirm the
guarantee on an operating LPG or petrochemical column with live distributed-control-system data.
